\section{Conclusion}
\label{sec:conclusion}

The seats-votes curve is the right exhibit to lead with in partisan fairness expert
testimony, and this paper has provided the analytical foundation for that recommendation.
The sufficiency result---EG and Partisan Bias are both recoverable from $S(v)$---means
that no scalar metric adds information beyond what the curve already contains. The
graphical format communicates partisan asymmetry directly to courts without requiring
expert testimony about wasted-vote algebra or distributional statistics. And the
curve's legal track record---admitted in NC and OH state court redistricting
proceedings post-\textit{Rucho}---demonstrates that it is judicially recognized
as appropriate evidence.

The responsiveness statistic ($R = dS/dv|_{v=0.50}$) adds to the curve's informative
value by quantifying what competes with Partisan Bias as the most important courtroom
argument: suppressed electoral accountability. When the enacted NC map produces $R = 1.3$
vs.\ \textsc{bisect}'s $R = 2.0$, the practical argument is: Republican incumbents
under the enacted map are insulated from vote swings that would cost them their seats
under a neutral map. This is not a proportionality claim (federally impermissible
post-\textit{Rucho}) but an accountability claim: the map suppresses the ability of
voters to change their representation through elections, which is precisely what
state constitutional ``free elections'' clauses are designed to protect.

The comparison of bisect and enacted seats-votes curves provides state courts with
the single most complete visual exhibit of partisan asymmetry available:
\begin{itemize}
\item One curve passes through $(0.50, 0.50)$ with slope $2.0$: neutral, responsive,
  symmetric.
\item The other curve passes through $(0.50, 0.43)$ with slope $1.3$: biased,
  unresponsive, asymmetric.
\item Both are drawn from the same state's geography; the difference reflects
  the mapmaker's deliberate choices, not geographic constraints.
\end{itemize}
No scalar metric captures this contrast as clearly, completely, or intuitively.

\bigskip

\noindent\textit{The author thanks the redistricting research team for data
access and algorithmic development. All errors are the author's own.}
